\documentclass{beamer}
\mode<presentation>
\usepackage{amsmath}
\usepackage{amssymb}
%\usepackage{advdate}
\usepackage{graphicx}
\graphicspath{{./figs/}}
\usepackage{adjustbox}
\usepackage{subcaption}
\usepackage{enumitem}
\usepackage{multicol}
\usepackage{mathtools}
\usepackage{listings}
\usepackage{url}
\def\UrlBreaks{\do\/\do-}
\usetheme{Boadilla}
\usecolortheme{lily}
\setbeamertemplate{footline}
{
  \leavevmode%
  \hbox{%
  \begin{beamercolorbox}[wd=\paperwidth,ht=2.25ex,dp=1ex,right]{author in head/foot}%
    \insertframenumber{} / \inserttotalframenumber\hspace*{2ex} 
  \end{beamercolorbox}}%
  \vskip0pt%
}
\setbeamertemplate{navigation symbols}{}

\providecommand{\nCr}[2]{\,^{#1}C_{#2}} % nCr
\providecommand{\nPr}[2]{\,^{#1}P_{#2}} % nPr
\providecommand{\mbf}{\mathbf}
\providecommand{\pr}[1]{\ensuremath{\Pr\left(#1\right)}}
\providecommand{\qfunc}[1]{\ensuremath{Q\left(#1\right)}}
\providecommand{\sbrak}[1]{\ensuremath{{}\left[#1\right]}}
\providecommand{\lsbrak}[1]{\ensuremath{{}\left[#1\right.}}
\providecommand{\rsbrak}[1]{\ensuremath{{}\left.#1\right]}}
\providecommand{\brak}[1]{\ensuremath{\left(#1\right)}}
\providecommand{\lbrak}[1]{\ensuremath{\left(#1\right.}}
\providecommand{\rbrak}[1]{\ensuremath{\left.#1\right)}}
\providecommand{\cbrak}[1]{\ensuremath{\left\{#1\right\}}}
\providecommand{\lcbrak}[1]{\ensuremath{\left\{#1\right.}}
\providecommand{\rcbrak}[1]{\ensuremath{\left.#1\right\}}}
\theoremstyle{remark}
\newtheorem{rem}{Remark}
\newcommand{\sgn}{\mathop{\mathrm{sgn}}}
\providecommand{\abs}[1]{\left\vert#1\right\vert}
\providecommand{\res}[1]{\Res\displaylimits_{#1}} 
\providecommand{\norm}[1]{\lVert#1\rVert}
\providecommand{\mtx}[1]{\mathbf{#1}}
\providecommand{\mean}[1]{E\left[ #1 \right]}
\providecommand{\fourier}{\overset{\mathcal{F}}{ \rightleftharpoons}}
%\providecommand{\hilbert}{\overset{\mathcal{H}}{ \rightleftharpoons}}
\providecommand{\system}[1]{\overset{\mathcal{#1}}{ \longleftrightarrow}}
%\providecommand{\system}{\overset{\mathcal{H}}{ \longleftrightarrow}}
	%\newcommand{\solution}[2]{\textbf{Solution:}{#1}}
%\newcommand{\solution}{\noindent \textbf{Solution: }}
\providecommand{\dec}[2]{\ensuremath{\overset{#1}{\underset{#2}{\gtrless}}}}
\newcommand{\myvec}[1]{\ensuremath{\begin{pmatrix}#1\end{pmatrix}}}
\let\vec\mathbf

\lstset{
%language=C,
frame=single, 
breaklines=true,
columns=fullflexible
}

\numberwithin{equation}{section}

\title{12.100}
\author{ai25btech11022 - Narshitha}
% \maketitle
% \newpage
% \bigskip
\begin{document}
{\let\newpage\relax\maketitle}
\renewcommand{\thefigure}{\theenumi}
\renewcommand{\thetable}{\theenumi} 
\textbf{Question}:\\
A Markov chain with state space \brak{0,1,2} has transition matrix \\
\begin{align}
    \myvec{\frac{1}{2} && \frac{1}{2} && 0\\\frac{1}{2} && \frac{1}{2} && 0 \\ \frac{1}{3}&& \frac{1}{3}&&\frac{1}{3}}
\end{align}
Which Statements are true?
\begin{enumerate}
    \item 0 and 1 are recurrent
    \item 2 is transient
    \item Chain has unique stationary distribution
    \item Chain is irreducible
\end{enumerate}
\textbf{Solution}:\\
\textbf{Markov Chain}:\\
The probability of going next state j  only depends upon current state i.
\begin{align}
    P\brak{X_{n+1} =j|X_n=i,X_{n-1},...,X_0}=P\brak{X_{n+1} =j|X_n=i}
\end{align}
\textbf{Transition Matrix}:
\begin{align}
\vec{P}=[ p_{ij}]\\
p_{ij}=P\brak{X_{n+1}=j|X_n =i}\\
\Sigma p_{ij} =1
\end{align}
\textbf{Recurrent}:Recurrent state is the one that the chain will return to with probability 1.\\
\textbf{Transient}: There is non zero probability that it will never return.\\
As \begin{align}
    p_{01} >0 ,p_{10}>0
\end{align} 0 and 1 are recurrent as they are closed communicating class.
As
\begin{align}
    p_{02}=0,p_{12}=0
\end{align} 2 is transient\\
\textbf{Stationary distribution}: Stationary distribution  is $\vec{\pi}$ such that
\begin{align}
    \vec{\pi}.\vec{P}=\vec{\pi} ,\Sigma \pi_i=1\\
    \implies  \vec{\pi}.\brak{\vec{P-I}}=0\\
    \brak{\vec{P^T-I}}.\vec{\pi}^T =0\\
    \myvec{1 &&1&&1\\\frac{-1}{2}&&\frac{1}{2}&& \frac{1}{3}\\\frac{1}{2} && \frac{-1}{2}&& \frac{1}{3}\\0 && 0&& \frac{-2}{3}}.\vec{\pi}^T =\myvec{1\\0\\0\\0}
\end{align}
The augmented matrix is
\begin{align}
    \myvec{1 &&1&&1 && \vrule && 1\\\frac{-1}{2}&&\frac{1}{2}&& \frac{1}{3} && \vrule && 0\\\frac{1}{2} && \frac{-1}{2}&& \frac{1}{3}&& \vrule && 0\\0 && 0&& \frac{-2}{3} && \vrule && 0}
\end{align}
$R_2 \longrightarrow R_2+R_3/2 $
\begin{align}
    \myvec{1 &&1&&1 && \vrule && 1\\\frac{-1}{2}&&\frac{1}{2}&& 0 && \vrule && 0\\\frac{1}{2} && \frac{-1}{2}&& \frac{1}{3}&& \vrule && 0\\0 && 0&& \frac{-2}{3} && \vrule && 0}\\
    \vec{{\pi}}^T=\myvec{\frac{1}{2}&& \frac{1}{2} && 0}
\end{align}
Therefore Chain has unique stationary distribution .\\
\textbf{Irreducible chain}: If for every j there exist $\brak{\vec{P}^n}_{ij}> 0$\\
the chain is not irreducible as 2 for j=2 ,The condition is not satisfied
\end{document}

